% Generated draft: requires mathematical and source review.
\documentclass[11pt]{article}
\usepackage[margin=25mm]{geometry}
\usepackage{fontspec}
\setmainfont{Latin Modern Roman}
\usepackage{amsmath,amssymb,mathtools,unicode-math}
\setmathfont{Latin Modern Math}
\usepackage{enumitem,needspace,xurl,hyperref}
\setlength{\parindent}{0pt}
\setlength{\parskip}{5pt}
\setlength{\emergencystretch}{3em}
\newcommand{\N}{\mathbb{N}}
\newcommand{\Z}{\mathbb{Z}}
\newcommand{\Q}{\mathbb{Q}}
\newcommand{\R}{\mathbb{R}}
\newcommand{\C}{\mathbb{C}}
\newcommand{\F}{\mathbb{F}}
\newcommand{\A}{\mathbb{A}}
\newcommand{\PP}{\mathbb{P}}
\newcommand{\E}{\mathbb{E}}
\newcommand{\eps}{\varepsilon}
\newcommand{\dd}{\,\mathrm{d}}
\begin{document}
\section*{UnsolvedMath ID: 20002728}
\textbf{Source problem code:} catalogue entry
\subsection{Definitions and mathematical statement}
A. Zelevinsky's Problem 6.5 asks for a tame/wild classification of infinite-type cluster algebras extending analogous notions for Artin groups and quivers. The split ``Problem 6. 5.'' is typographical. The source remark says affine types should certainly be tame, but no general definitions, coefficient conventions, equivalence relation or classification criterion are specified. The aid tests candidate invariants in coefficient-free rank two rather than supplying a general tame/wild theory.

Let $b,c$ be positive integers and let $\mathcal A(b,c)$ have initial principal exchange matrix
\[
B=\begin{pmatrix}0&b\\-c&0\end{pmatrix}.
\]
Its cluster variables satisfy rank-two exchange relations $x_{j-1}x_{j+1}=1+x_j^{d_j}$ with exponents alternating between $b$ and $c$, according to the indexing convention. The labeled mutation class of the $2\times2$ principal part is $\{B,-B\}$; this is not a statement about an extended principal-coefficient matrix. For every infinite case $bc\geq4$, the exchange graph is a bi-infinite path, and a radius-$R$ ball contains $2\lfloor R\rfloor+1$ clusters. Thus mutation-class cardinality and exchange-ball growth degree cannot distinguish affine $bc=4$ from indefinite $bc>4$.

The aid instead uses denominator growth. For a Laurent cluster variable in the initial variables, its denominator vector records the two denominator exponents after cancelling common monomial factors. Along a fixed parity of an infinite Coxeter ray, index $v_j$ by successive two-step Coxeter moves. The supplied growth statements are
\[
\|v_j\|=\Theta(j)\quad(bc=4),\qquad
\|v_j\|=\Theta(\Lambda^j)\quad(bc>4),
\]
\[
\Lambda=\frac{bc-2+\sqrt{(bc-2)^2-4}}2>1,
\qquad
\lim_{j\to\infty}\frac1j\log\max\{1,\|v_j\|\}=\log\Lambda\quad(bc>4).
\]
The norm can be any fixed norm in dimension two, and finite initial terms do not affect the rate. The entropy is per \emph{two-step} Coxeter move, not per single mutation. In finite cases $bc\leq3$, denominator vectors are bounded and periodic; hence the rank-two growth trichotomy is bounded/periodic, linear, and exponential. Entropy alone is zero in both the finite and affine cases.

For the skew-symmetric Kronecker subfamily $b=c=m$, this agrees with the classical representation-type boundary: the two-arrow Kronecker quiver is tame, while the $m$-arrow quiver for $m\geq3$ is wild, in the standard algebraically closed-field representation setting. Therefore any proposed cluster invariant intended to separate those examples cannot factor solely through principal mutation-class size or exchange-growth degree. The denominator rate supplies a separator for these rank-two tests; it does not prove that all zero-entropy cluster algebras are tame, all positive-entropy algebras are wild, or that the same definition extends to every skew-symmetrizable or Artin setting.
\subsection{Short English statement}
The source seeks a general tame/wild theory but leaves the terms undefined. The aid shows why two simple rank-two invariants are insufficient: every infinite rank-two exchange graph is a path and every principal mutation class has the same size. Denominator growth separates affine from indefinite rank two and matches the Kronecker examples, without providing the requested global classification.
\subsection{Sources}
\begin{itemize}
\item Local UnsolvedMath registry, permanent ID $20002728$, \texttt{catalogue entry}; \texttt{aim-probability-notes.json}, \texttt{notes[169]}, extracted source item 6.
\item American Institute of Mathematics, \emph{Braid groups, clusters and free probability}, Problem 6.5 (A. Zelevinsky), including ``Affine types are certainly tame,'' \url{https://aimath.org/WWN/braidgroups/braidgroups.pdf}. Supplied verification identifies the split heading as an extraction artifact.
\item Research attempt 1: partial result, problem partially solved at run, rigor label ``proved,'' novelty confidence low. Registry classification partial progress, status partially solved, statement status exact. No authoritative separate references for the rank-two growth formulas or representation-type comparison are supplied.
\end{itemize}
\end{document}
